\section{Ethics}
\label{sec:ethics}

This work involves no human subjects. 
The RQ2 probes are the unauthenticated handshakes a browser starts. Each
endpoint received the probe's ClientHello and one standard handshake, each
retried a bounded number of times on failure, and no application data. The global rate was
\data{rqtwo-rate} connections per second and at most one connection per host
was open at a time, both enforced in code. Nothing was exploited and no
credentials were tried. The project's documentation describes exactly what a
probe sends and how to opt out; the opt-out list is honoured by every run and
held \data{rqtwo-exclude} entries at the time of the scan. We report results
in aggregate and will report notable findings about federal systems through
CISA. The scanning source address has no reverse DNS record pointing to the
project, which we note as a shortcoming of this run. RQ3 clones public
repositories and runs none of their code. These practices follow the
recommendations of the ZMap authors~\cite{zmap} and the Menlo
Report~\cite{menlo}.
